% !TeX program = xelatex
% !TeX TXS-program:compile = txs:///xelatex
% preamble-exam-zh.tex — 极简讲解页共享 preamble
% 目标：复刻「题目独立、提示轻框、路线箭头、主体双栏、答案轻收束」的讲义风格。

\documentclass{exam-zh}

% ---------- 基础配置 ----------
\examsetup{
  page/size = a4paper,
  page/show-head = false,
  page/show-foot = false,
  sealline/show = false,
  font = none,
  math-font = none,
}

% ---------- 包 ----------
\usepackage{geometry}
\geometry{top=10mm,bottom=14mm,left=12mm,right=10mm}
\AtBeginDocument{\newgeometry{top=10mm,bottom=14mm,left=12mm,right=10mm}}

\usepackage{xcolor}
\usepackage{needspace}
\usepackage{enumitem}
\usepackage{array}
\usepackage{tabularx}
\usepackage{tcolorbox}
\tcbuselibrary{skins,breakable}
\usepackage{paracol}
\usepackage{graphicx}
\usepackage{tikz}
\usetikzlibrary{calc,intersections,angles,quotes,arrows.meta,decorations.markings,patterns.meta,backgrounds}
\usepackage{pgfplots}
\pgfplotsset{compat=1.18}

% ---------- 打印/屏幕颜色切换 ----------
% 用法：\documentclass[monochrome]{exam-zh} 可降级为灰度。
% 注意：不要使用 \if@xxx 放在 makeatother 外部；这里改成普通条件名。
\newif\ifeduprintcolor
\eduprintcolortrue
\makeatletter
\@ifclasswith{exam-zh}{monochrome}{\eduprintcolorfalse}{}
\makeatother

% ---------- 语义颜色 ----------
\ifeduprintcolor
  \definecolor{edu-blue}{RGB}{31,62,126}
  \definecolor{edu-blue-soft}{RGB}{232,238,250}
  \definecolor{edu-blue-bg}{RGB}{248,250,254}
  \definecolor{edu-ink}{RGB}{30,35,45}
  \definecolor{edu-muted}{RGB}{118,124,136}
  \definecolor{edu-line}{RGB}{209,218,235}
  \definecolor{edu-soft-bg}{RGB}{250,251,253}
  \definecolor{edu-red}{RGB}{168,58,58}
  \definecolor{edu-red-bg}{RGB}{255,247,245}
\else
  \definecolor{edu-blue}{gray}{0.15}
  \definecolor{edu-blue-soft}{gray}{0.90}
  \definecolor{edu-blue-bg}{gray}{0.98}
  \definecolor{edu-ink}{gray}{0.10}
  \definecolor{edu-muted}{gray}{0.45}
  \definecolor{edu-line}{gray}{0.82}
  \definecolor{edu-soft-bg}{gray}{0.97}
  \definecolor{edu-red}{gray}{0.28}
  \definecolor{edu-red-bg}{gray}{0.97}
\fi
\colorlet{edu-diagram-muted}{edu-muted}

% ---------- TikZ diagram semantic macros ----------
\definecolor{edu-diagram-ink}{RGB}{17,24,39}
\definecolor{edu-diagram-marker}{RGB}{220,38,38}
\definecolor{edu-diagram-angle}{RGB}{5,150,105}
\tikzset{
  diagram point/.style={fill=edu-diagram-ink},
  diagram tick/.style={draw=edu-diagram-marker,line width=0.9pt,line cap=round},
  diagram segment/.style={draw=edu-diagram-ink,line cap=round},
  point label/.style={inner sep=1pt},
  condition label/.style={inner sep=1pt}
}
\providecommand{\DrawSegment}[3][]{\draw[diagram segment,#1] (#2) -- (#3);}
\providecommand{\DrawDashedSegment}[3][]{\draw[diagram segment,dashed,#1] (#2) -- (#3);}
\providecommand{\Triangle}[4][]{\path[#1] (#2) -- (#3) -- (#4) -- cycle;}
\providecommand{\Quadrilateral}[5][]{\path[#1] (#2) -- (#3) -- (#4) -- (#5) -- cycle;}
\providecommand{\PolygonPath}[2][]{\path[#1] #2 -- cycle;}
\providecommand{\DiagramPointRadius}{0.052cm}
\providecommand{\PointDot}[2][]{\fill[diagram point,#1] (#2) circle[radius=\DiagramPointRadius];}
\providecommand{\PointLabel}[3][]{\node[point label,#1] at (#2) {#3};}
\providecommand{\SegmentLabel}[4][]{\node[condition label,#1] at ($(#2)!0.5!(#3)$) {#4};}
\providecommand{\CoordinateTag}[3][]{\node[condition label,#1] at (#2) {#3};}
\providecommand{\AngleMark}[4][]{\pic[draw=edu-diagram-angle,line width=0.9pt,angle radius=0.48cm,#1] {angle=#2--#3--#4};}
\providecommand{\AngleLabel}[5][]{\pic["{#5}",draw=none,angle radius=0.64cm,#1] {angle=#2--#3--#4};}
\providecommand{\RightAngleMark}[4][]{\pic[draw=edu-diagram-marker,line width=0.9pt,angle radius=0.28cm,#1] {right angle=#2--#3--#4};}
\providecommand{\EqualTick}[3][]{%
  \path[postaction={decorate},decoration={markings,mark=at position .5 with {\draw[diagram tick,#1] (0pt,-4pt) -- (0pt,4pt);}}] (#2) -- (#3);%
}
\providecommand{\DoubleEqualTick}[3][]{%
  \path[postaction={decorate},decoration={markings,mark=at position .46 with {\draw[diagram tick,#1] (0pt,-4pt) -- (0pt,4pt);},mark=at position .54 with {\draw[diagram tick,#1] (0pt,-4pt) -- (0pt,4pt);}}] (#2) -- (#3);%
}
\providecommand{\TripleEqualTick}[3][]{%
  \path[postaction={decorate},decoration={markings,mark=at position .43 with {\draw[diagram tick,#1] (0pt,-4pt) -- (0pt,4pt);},mark=at position .5 with {\draw[diagram tick,#1] (0pt,-4pt) -- (0pt,4pt);},mark=at position .57 with {\draw[diagram tick,#1] (0pt,-4pt) -- (0pt,4pt);}}] (#2) -- (#3);%
}
\providecommand{\ParallelMark}[3][]{\EqualTick[#1]{#2}{#3}}
\providecommand{\NamedSegmentPath}[3]{\path[name path=#1] (#2) -- (#3);}
\providecommand{\IntersectPaths}[3]{\path[name intersections={of=#2 and #3, by=#1}];}

% ---------- 全局排版 ----------
\setlength{\parindent}{0pt}
\setlength{\parskip}{0pt}
\linespread{1.16}
\renewcommand{\arraystretch}{1.15}

% ctex 通常提供 \kaishu；这里用它做教师旁白感。
\newcommand{\edukai}{\kaishu}

% ---------- 顶部元信息 ----------
\newcommand{\eduTopBar}[2]{%
  \noindent
  \begin{tabularx}{\linewidth}{@{}Xr@{}}
    {\color{edu-blue}\bfseries #1} &
    {\color{edu-blue}\bfseries #2}
  \end{tabularx}
  \vspace{8mm}
}

% ---------- 轻标题体系 ----------
% 只有真正的大结构才用它；不要给提示/路线滥用标题。
\newcommand{\eduSectionTitle}[1]{%
  \needspace{5\baselineskip}%
  \vspace{2mm}
  {\color{edu-blue}\bfseries\large #1}\par
  \vspace{3mm}
}

\newcommand{\eduSubTitle}[1]{%
  \needspace{4\baselineskip}%
  \vspace{1mm}
  {\color{edu-blue}\bfseries #1}\par
  \vspace{1mm}
}

% ---------- 小问题干：复用 exam (1)(2)(3) 格式 ----------
\newenvironment{eduSubQuestionItem}[1]{%
  \needspace{4\baselineskip}%
  \vspace{2mm}
  \par\noindent
  \hangindent=8mm\hangafter=1
  {\color{edu-blue}\bfseries #1}\enspace
  \begingroup
  \color{edu-ink}
}{%
  \endgroup
  \par
  \vspace{3mm}
}

% ---------- 辅助信息条（解题流程/关键想法/思考）楷体 ----------
\newenvironment{eduAuxItem}[1]{%
  \needspace{4\baselineskip}%
  \vspace{2mm}
  \par\noindent
  \hangindent=8mm\hangafter=1
  {\color{edu-blue}\edukai $\triangleright$~#1}\enspace
  \begingroup
  \color{edu-ink}
  \edukai
}{%
  \endgroup
  \par
  \vspace{4mm}
}

\newcommand{\eduColumnTitle}[2]{%
  {\color{edu-blue}\bfseries #1\hspace{1.5mm}#2}\par
  \vspace{2.5mm}
}

% ---------- 图标：不用外部 icon 字体，保证可移植 ----------
\newcommand{\eduIconBulb}{\(\circ\)}
\newcommand{\eduIconBook}{\(\square\)}
\newcommand{\eduIconCheck}{\(\checkmark\)}
\newcommand{\eduIconPin}{\(\diamond\)}
\newcommand{\eduIconNote}{\(\triangleright\)}

% ---------- 题目区 ----------
% 题目区是页面入口：弱标签 + 一条长蓝线 + 大题干。
\newenvironment{eduProblemBlock}[1][题目]{%
  \needspace{8\baselineskip}
  {\small\color{edu-muted}#1}\par
  \vspace{1.5mm}
  {\color{edu-blue}\hrule height 0.55pt}
  \vspace{4.5mm}
  \begingroup
  \color{edu-ink}
  \large
  \setlength{\parskip}{2.2mm}
  \linespread{1.20}\selectfont
}{%
  \par
  \endgroup
  \vspace{6mm}
}

\newenvironment{eduSubQuestions}{%
  \begin{enumerate}[
    label=(\arabic*),
    leftmargin=8mm,
    itemsep=2.0mm,
    topsep=2.0mm,
    parsep=0pt
  ]
}{%
  \end{enumerate}
}

% ---------- 读题提示：轻框，不做同级标题 ----------
\newtcolorbox{eduReadTipBox}{
  enhanced,
  breakable,
  colback=white,
  colframe=edu-blue!45,
  boxrule=0.45pt,
  arc=1.6mm,
  left=10mm,
  right=4mm,
  top=5mm,
  bottom=5mm,
  before skip=1mm,
  after skip=7mm,
  fontupper=\edukai\normalsize,
  colupper=edu-ink,
  overlay={
    \node[anchor=north west, text=edu-blue] at ([xshift=3.2mm,yshift=-3.2mm]frame.north west)
    {\small\eduIconNote};
  }
}

\newenvironment{eduTipList}{%
  \begin{itemize}[
    label=\textbullet,
    leftmargin=5mm,
    itemsep=1.2mm,
    topsep=0pt,
    parsep=0pt
  ]
}{%
  \end{itemize}
}

% ---------- 解题路线：虚线框 + 文字箭头 ----------
\newenvironment{eduRouteFlow}{%
  \needspace{4\baselineskip}%
  \vspace{1mm}
  \begin{center}
  \normalsize
}{%
  \end{center}
  \vspace{7mm}
}
\newcommand{\eduRouteStep}[1]{%
  {\color{edu-ink}#1}%
}
\newcommand{\eduRouteArrow}{%
  \;$\boldsymbol{\to}$\;%
}

% ---------- 主体讲解：使用 paracol 实现跨页自适应双栏 ----------
\newenvironment{eduDualExplain}{%
  \needspace{6\baselineskip}% 防止标题孤行
  \columnratio{0.26, 0.74}%
  \setlength{\columnsep}{8mm}%
  \columnseprule=0.45pt%
  \colseprulecolor{edu-blue}%
  \begin{paracol}{2}% 开启双栏
}{%
  \end{paracol}% 结束双栏
  \vspace{5mm}
}

\newenvironment{eduExplainLeft}{%
  \setlength{\linewidth}{\columnwidth}%
  \hsize=\columnwidth
  \normalsize\color{edu-ink}%
}{%
  \switchcolumn
}

\newcommand{\eduColumnDivider}{}

\newenvironment{eduExplainRight}{%
  \setlength{\linewidth}{\columnwidth}%
  \hsize=\columnwidth
  \normalsize\color{edu-ink}%
}{%
}

\newtcolorbox{eduSideHintBox}[1]{
  enhanced,
  breakable,
  colback=edu-blue-bg,
  colframe=edu-line,
  boxrule=0.35pt,
  borderline west={1.5pt}{0pt}{edu-blue},
  arc=0.8mm,
  left=2.4mm,
  right=2.4mm,
  top=2mm,
  bottom=2mm,
  before skip=1mm,
  after skip=1.5mm,
  title={\color{edu-blue}\bfseries\small #1},
  coltitle=edu-blue,
  attach title to upper={\par\vspace{0.7mm}},
  fontupper=\small
}

\newtcolorbox{eduSideMistakeBox}[1]{
  enhanced,
  breakable,
  colback=edu-red-bg,
  colframe=edu-red!35,
  boxrule=0.35pt,
  borderline west={1.5pt}{0pt}{edu-red},
  arc=0.8mm,
  left=2.4mm,
  right=2.4mm,
  top=2mm,
  bottom=2mm,
  before skip=1mm,
  after skip=1.5mm,
  title={\color{edu-red}\bfseries\small #1},
  coltitle=edu-red,
  attach title to upper={\par\vspace{0.7mm}},
  fontupper=\small
}

\newtcolorbox{eduSideNoteBox}[1]{
  enhanced,
  breakable,
  colback=edu-soft-bg,
  colframe=edu-line,
  boxrule=0.35pt,
  borderline west={1.5pt}{0pt}{edu-muted},
  arc=0.8mm,
  left=2.4mm,
  right=2.4mm,
  top=2mm,
  bottom=2mm,
  before skip=1mm,
  after skip=1.5mm,
  title={\color{edu-muted}\bfseries\small #1},
  coltitle=edu-muted,
  attach title to upper={\par\vspace{0.7mm}},
  fontupper=\small
}

\newcommand{\eduSideHint}[2]{%
  \begin{eduSideHintBox}{#1}#2\end{eduSideHintBox}%
}
\newcommand{\eduSideMistake}[2]{%
  \begin{eduSideMistakeBox}{#1}#2\end{eduSideMistakeBox}%
}
\newcommand{\eduSideNote}[2]{%
  \begin{eduSideNoteBox}{#1}#2\end{eduSideNoteBox}%
}

\newcommand{\eduExplainStep}[3]{%
  \needspace{3\baselineskip}%
  \vspace{1.2mm}%
  {\color{edu-blue}\bfseries step#1.\enspace #2}\par
  \vspace{0.8mm}%
  \begingroup
  \color{edu-ink}#3\par
  \endgroup
  \vspace{2.2mm}%
}

\newenvironment{eduNumberList}{%
  \begin{enumerate}[
    label=\arabic*.,
    leftmargin=6mm,
    itemsep=4.8mm,
    topsep=0pt,
    parsep=0pt
  ]
}{%
  \end{enumerate}
}

\newenvironment{eduBulletList}{%
  \begin{itemize}[
    label=\textbullet,
    leftmargin=5mm,
    itemsep=1.2mm,
    topsep=0pt,
    parsep=0pt
  ]
}{%
  \end{itemize}
}

% ---------- 底部双栏：答案 + 方法提醒 ----------
\newenvironment{eduBottomGrid}{%
  \columnratio{0.32, 0.68}%
  \setlength{\columnsep}{11mm}%
  \columnseprule=0pt%
  \begin{paracol}{2}%
}{%
  \end{paracol}%
}

\newenvironment{eduSummaryLeft}{%
}{%
}

\newcommand{\eduSummaryDivider}{%
  \switchcolumn
}

\newenvironment{eduSummaryRight}{%
}{%
}

% ---------- 答案/方法提醒：轻量收束 ----------
\newtcolorbox{eduSummaryBlock}[2]{
  enhanced,
  breakable,
  colback=edu-blue-bg,
  colframe=edu-line,
  boxrule=0.45pt,
  arc=1.2mm,
  left=3.5mm,
  right=3.5mm,
  top=3mm,
  bottom=3mm,
  before skip=1mm,
  after skip=2.5mm,
  title={\color{edu-blue}\bfseries #1\hspace{1.5mm}#2},
  coltitle=edu-blue,
  attach title to upper={\par\vspace{1.5mm}},
  fontupper=\normalsize
}

% ---------- 例题单栏解答卡片 ----------
\newtcolorbox{eduSolutionBox}[1]{
  enhanced,
  breakable,
  colback=white,
  colframe=edu-blue!40,
  boxrule=0.55pt,
  arc=1.2mm,
  left=4mm,
  right=4mm,
  top=3.5mm,
  bottom=3.5mm,
  before skip=2mm,
  after skip=3mm,
  title={\color{edu-blue}\bfseries \eduIconBook\hspace{1.5mm}#1},
  coltitle=edu-blue,
  attach title to upper={\par\vspace{1.5mm}},
  fontupper=\normalsize
}

% ---------- 变式训练：完整题目 + 学生留白 ----------
\newtcolorbox{eduVariationTraining}[1]{
  enhanced,
  breakable,
  colback=white,
  colframe=edu-blue!45,
  boxrule=0.5pt,
  borderline west={1.8pt}{0pt}{edu-blue},
  arc=1.2mm,
  left=4mm,
  right=4mm,
  top=3.2mm,
  bottom=3.2mm,
  before skip=2mm,
  after skip=4mm,
  title={\color{edu-blue}\bfseries #1},
  coltitle=edu-blue,
  attach title to upper={\par\vspace{1.8mm}},
  fontupper=\normalsize,
  colupper=edu-ink
}

\newenvironment{eduPlainList}{%
  \begin{enumerate}[
    label=(\arabic*),
    leftmargin=7mm,
    itemsep=1.2mm,
    topsep=0pt,
    parsep=0pt
  ]
}{%
  \end{enumerate}
}

% ---------- 兼容旧 block 的轻量盒子 ----------
\newtcolorbox{eduSoftBox}{
  enhanced,
  breakable,
  colback=edu-soft-bg,
  colframe=edu-line,
  boxrule=0.35pt,
  arc=1mm,
  left=3mm,
  right=3mm,
  top=2mm,
  bottom=2mm,
  before skip=1mm,
  after skip=2mm,
  fontupper=\small
}

\newtcolorbox{eduMistakeBox}{
  enhanced,
  breakable,
  colback=edu-red-bg,
  colframe=edu-red!40,
  boxrule=0.35pt,
  arc=1mm,
  left=3mm,
  right=3mm,
  top=2.5mm,
  bottom=2.5mm,
  before skip=1mm,
  after skip=2mm,
  fontupper=\normalsize
}

\newtcolorbox{eduLegacySolution}{
  enhanced,
  breakable,
  colback=edu-soft-bg,
  colframe=edu-line,
  boxrule=0.35pt,
  arc=1mm,
  left=3mm,
  right=3mm,
  top=2mm,
  bottom=2mm,
  before skip=-2mm,
  after skip=4mm,
  fontupper=\small
}

\newtcolorbox{eduTeacherNote}{
  enhanced,
  breakable,
  colback=edu-soft-bg,
  colframe=edu-line,
  boxrule=0pt,
  borderline west={1.5pt}{0pt}{edu-muted},
  left=2.5mm,
  right=2.5mm,
  top=2mm,
  bottom=2mm,
  before skip=1mm,
  after skip=2mm,
  fontupper=\footnotesize,
  colupper=edu-muted
}

% ---------- 答题区 ----------
\newcommand{\answerarea}[1][40mm]{%
  \vspace{3mm}%
  \begin{tcolorbox}[
    enhanced,
    breakable,
    colback=white,
    colframe=edu-line,
    boxrule=0.55pt,
    arc=0.8mm,
    height=#1,
    valign=top,
  ]
  \end{tcolorbox}%
}

\newsavebox{\diagramtikzbox}

\newenvironment{diagramcoltikz}[2]{%
  \begin{minipage}[t]{#1}
    \vspace{0pt}\centering
    \def\diagramcaption{#2}%
    \begin{lrbox}{\diagramtikzbox}%
}{%
    \end{lrbox}%
    \usebox{\diagramtikzbox}%
    \ifx\diagramcaption\empty\else
      \par\vspace{1mm}{\scriptsize\color{edu-diagram-muted}\diagramcaption}
    \fi
  \end{minipage}%
}

\newenvironment{diagramrowtikz}[3]{%
  \begin{minipage}[t]{#1}
    \vspace{0pt}\centering
    \def\diagramlabel{#2}%
    \def\diagramcaption{#3}%
    \begin{lrbox}{\diagramtikzbox}%
}{%
    \end{lrbox}%
    \usebox{\diagramtikzbox}%
    \ifx\diagramlabel\empty\else
      \par\vspace{1mm}{\scriptsize\bfseries\diagramlabel}
    \fi
    \ifx\diagramcaption\empty\else
      \par\vspace{0.5mm}{\scriptsize\color{edu-diagram-muted}\diagramcaption}
    \fi
  \end{minipage}%
}

\newenvironment{answerareawithdiagramtikz}[3]{%
  \par\vspace{3mm}%
  \noindent
  \begin{minipage}[t]{\dimexpr\linewidth-#2-6mm\relax}
    \vspace{0pt}\rule{0pt}{#1}
  \end{minipage}\hfill
  \begin{diagramcoltikz}{#2}{#3}
}{%
  \end{diagramcoltikz}
  \par\vspace{2mm}%
}
